\documentclass[11pt]{article}
\usepackage[a4paper,margin=22mm]{geometry}
\usepackage{amsmath,amssymb}
\usepackage{graphicx}
\usepackage{microtype}
\usepackage[hidelinks]{hyperref}

\graphicspath{{../../results/collective_agent_statmech_v15/figures/pdf/}}
\newcommand{\resultfigure}[3]{%
  \begin{figure}[p]
  \centering
  \IfFileExists{../../results/collective_agent_statmech_v15/figures/pdf/#1}{%
    \includegraphics[width=#2]{#1}%
  }{\fbox{\parbox[c][55mm][c]{0.88\textwidth}{\centering
    Figure generated by the frozen V15 reproduction workflow.}}}%
  \caption{#3}
  \end{figure}%
}

\title{Supplementary material for\\
Statistical-mechanical characterization of memory and quench response in
state-separated LLM-agent networks}
\author{}
\date{}

\begin{document}
\maketitle

\section*{S1. Finite-window persistence and order-parameter shape}

The diagnostics below accompany equations (8)--(9) of the main text.  Each
quantity is computed within one complete model--graph--environment trajectory
before uncertainty is summarized across the six clusters per model.  The
autocorrelation sum uses the fixed two-sweep truncation, while one- and
three-sweep truncations remain machine-readable sensitivity analyses.  The
Binder statistic describes the shape of a finite $N=16$ magnetization
distribution.  The source-data table additionally retains full-phase versus
early/late half-phase estimates and cluster-mean versus pooled-moment rules.
It is not a crossing analysis or evidence of a phase transition.

\resultfigure{figure14_persistence_and_binder.pdf}{0.98\textwidth}{%
Finite-window persistence and order-parameter shape.  Autocorrelation curves
and their 95\% intervals compare the nominal late window with the three field
arms using complete clusters; integrated correlation sums show all cluster
values at the fixed two-sweep truncation.  Binder cumulants are paired with
empirical belief-magnetization occupancies so their finite-sample meaning
remains visible.  These diagnostics do not establish critical slowing down or
a phase transition.}

\end{document}
